\subsection{Faltings 半稳定性判据}\label{src:9.6.2-subsection}
我们建立 Faltings 半稳定性判据, 并应用它得到 $L_V$ 的半丰沛性.
\begin{theorem}[Faltings 半稳定性判据]\label{src:9.6.4}
设 $C$ 是代数闭域 $\Bbbk$ 上亏格 $g\geq2$ 的光滑连通射影曲线, $F$ 是 $C$ 上秩 $r$, 次数 $d$ 的向量丛. 置 $r'=r/\gcd(r,d)$, $d'=d/\gcd(r,d)$. 则下列条件等价:
\begin{enumerate}[(1)]
\item $F$ 半稳定;
\item 存在 $C$ 上向量丛 $V$, 满足 $\mathrm H^0(C,F\otimes V)=\mathrm H^1(C,F\otimes V)=0$;
\item 对 $m\gg0$ 及 $C$ 上秩 $mr'$, 次数 $mr'(g-1)-md'$ 的一般向量丛 $V$, 有 $\mathrm H^0(C,F\otimes V)=0$.
\end{enumerate}
\end{theorem}
\begin{remark}\label{src:9.6.5}
满足 (2) 或 (3) 的任意向量丛 $V$ 必半稳定; 交换 $F$ 与 $V$ 的角色, 应用 (2), (3) $\Rightarrow$ (1) 即得.
\end{remark}
\begin{proof}
(2) 中对 $V$ 的秩与次数的条件蕴含 $\chi(C,F\otimes V)=0$ (见注~\bookref{9.4.19}), 故 $\mathrm H^0(C,F\otimes V)=0$ 蕴含 $\mathrm H^1(C,F\otimes V)=0$. 这给出 (3) $\Rightarrow$ (2).

为说明 (2) $\Rightarrow$ (1), 应用 Riemann--Roch 得到
\begin{equation}\label{src:9.6.6}
\mu(F)+\mu(V)=\mu(F\otimes V)=\frac{\chi(F\otimes V)}{\mathrm{rk}(F\otimes V)}+g-1.
\end{equation}
假设 $F$ 不半稳定, $E\subseteq F$ 是使之不稳定的子层. 因为~\bookeqref{9.6.6} 对 $E$ 也成立, $\mu(E)>\mu(F)$ 蕴含 $\mu(E\otimes V)>\mu(F\otimes V)$, 从而
\[
\chi(E\otimes V)>\frac{\mathrm{rk}(E\otimes V)}{\mathrm{rk}(F\otimes V)}\chi(F\otimes V)=0.
\]
这蕴含 $\mathrm H^0(C,E\otimes V)\neq0$, 又因 $E\otimes V\subseteq F\otimes V$, 进一步蕴含 $\mathrm H^0(C,F\otimes V)\neq0$.

对较难的方向 (1) $\Rightarrow$ (3), $F$ 稳定时可直接由引理~\bookref{9.6.1}(2) 取 $\epsilon=0$ 得到. 为归结到稳定情形, 考虑伴随分次 $[\mathrm{gr}\,F]\in\cBun_{r,d}^{\mathrm{ss}}(\Bbbk)$, 它在 $[F]$ 的闭包中. $\mathrm H^0(C,F\otimes V)\neq0$ 的轨迹闭 (或等价地, 根据引理~\bookref{9.4.22}, 截面 $s_V$ 消失的轨迹闭). 所以若 $\mathrm H^0(C,\mathrm{gr}(F)\otimes V)=0$, 则 $\mathrm H^0(C,F\otimes V)=0$. 这将证明归结到 $F$ 多重稳定的情形. 将 $F=\bigoplus_iF_i$ 写为相同斜率的稳定丛的直和, 则对一般向量丛 $V$, $\mathrm H^0(C,F_i\otimes V)=0$ 蕴含 $\mathrm H^0(C,F\otimes V)=0$.

另见 \cite[定理 I.2, 第 514 页]{Fal93} (其中对 Higgs 丛给出更一般的论证), \cite[引理 3.1]{Ses93}, \cite[定理 2.4]{LP96}, \cite[定理 2.5]{Hei97} 及 \cite[命题 5.2]{ABB+22}.
\end{proof}
回忆特征零时, 已证明良模空间 $\cBun_{r,\xi}^{\mathrm{ss}}\to M_{r,\xi}^{\mathrm{ps}}$ 存在, 其中 $M_{r,\xi}^{\mathrm{ps}}$ 为固有代数空间 (定理~\bookref{9.3.11}). $\cBun_{r,\xi}^{\mathrm{ss}}$ 上行列式线丛 $\mathcal L_V$ 仅依赖于 $V$ 的秩与次数 (习题~\bookref{9.4.17}(2)). 因此取 $\mathrm{rk}\,V=r/\gcd(r,d)$, $\deg V=(r(g-1)-d)/\gcd(r,d)$, 就定义了 $\cBun_{r,\xi}^{\mathrm{ss}}$ 上良定义的线丛 $\mathcal L_V$. 根据命题~\bookref{9.4.29}, 此线丛下降为 $M_{r,\xi}^{\mathrm{ps}}$ 上的线丛 $L_V$.
\begin{corollary}\label{src:9.6.7}
设 $C$ 是特征零代数闭域 $\Bbbk$ 上亏格 $g\geq2$ 的光滑连通射影曲线. 对秩 $r/\gcd(r,d)$, 次数 $(r(g-1)-d)/\gcd(r,d)$ 的向量丛 $V$, 线丛 $L_V$ 在 $M_{r,\xi}^{\mathrm{ps}}$ 上半丰沛.
\end{corollary}
\begin{proof}
置 $r'=r/\gcd(r,d)$, $d'=d/\gcd(r,d)$. 根据 Faltings 半稳定性判据 (\bookref{9.6.4}), 对每个半稳定向量丛 $F$, 存在整数 $m$ 与秩 $mr'$, 次数 $mr'(g-1)-md'$ 的向量丛 $V'$, 使 $\mathrm H^0(C,F\otimes V')=0$. 由引理~\bookref{9.4.22}, 得到 $\mathcal L_{V'}\cong\mathcal L_V^{\otimes m}$ 的截面 $s_{V'}$, 满足 $s_{V'}([F])\neq0$. 因为
\[
\mathrm H^0(\cBun_{r,\xi}^{\mathrm{ss}},\mathcal L_V^{\otimes m})=\mathrm H^0(M_{r,\xi}^{\mathrm{ps}},L_V^{\otimes m}),
\]
这也给出 $\mathcal L_V^{\otimes m}$ 的截面 $s_{V'}$, 在 $[F]$ 的像处不消失.
\end{proof}
